% =============================================================================
% 2.9 Planification et synthèse de la méthodologie
% =============================================================================
\section{Planification et synthèse}
\label{sec:ch2-planification}

\subsection{Phases de développement}

Le développement s'organise en cinq phases, représentées à la
figure~\ref{fig:phases}. Les quatre premières constituent le périmètre de
la première version. La cinquième est optionnelle et conditionnée par les
résultats des précédentes.

\begin{figure}[htbp]
  \centering
  \resizebox{\linewidth}{!}{% Phases de développement
\begin{tikzpicture}[x=1.05cm, y=1cm]

  % axe des semaines
  \draw[btplight, line width=0.8pt] (0,0) -- (10.4,0);
  \foreach \x/\lab in {0/S1, 2/S3, 4/S5, 6/S7, 8/S9, 10/S11} {
    \draw[btplight, line width=0.8pt] (\x,0) -- (\x,-0.14);
    \node[btplabel, below=0pt] at (\x,-0.14) {\lab};
  }

  \newcommand{\phase}[5]{%
    \fill[#5, rounded corners=2pt] (#1,#3-0.24) rectangle (#2,#3+0.24);
    \node[font=\scriptsize\bfseries, text=white, anchor=west] at (#1+0.12,#3) {#4};
  }

  \phase{0}{2}{3.2}{Phase 1 — Socle}{btpdeep}
  \phase{2}{4}{2.5}{Phase 2 — Ingestion}{btpblue}
  \phase{4}{6}{1.8}{Phase 3 — Recherche hybride}{btpblue!85}
  \phase{6}{8}{1.1}{Phase 4 — Workflow et supervision}{btpsky}
  \phase{8}{10}{0.4}{Phase 5 — Graphe (option)}{btpslate}

  \node[btplabel, anchor=west] at (10.5,3.2) {Infrastructure, schéma, sécurité};
  \node[btplabel, anchor=west] at (10.5,2.5) {Analyse, découpage, indexation};
  \node[btplabel, anchor=west] at (10.5,1.8) {Fusion, ré-ordonnancement};
  \node[btplabel, anchor=west] at (10.5,1.1) {Génération, revue, export};
  \node[btplabel, anchor=west] at (10.5,0.4) {Entités et relations};

\end{tikzpicture}
}
  \caption{Phasage du développement}
  \label{fig:phases}
\end{figure}

\begin{description}
  \item[Phase 1 --- Socle.] Infrastructure, schéma relationnel,
        authentification, points d'entrée minimaux.
  \item[Phase 2 --- Ingestion.] Analyse structurelle, reconnaissance
        optique, découpage parent-enfant, indexation.
  \item[Phase 3 --- Recherche hybride.] BM25, vecteurs denses, fusion,
        ré-ordonnancement, construction de contexte.
  \item[Phase 4 --- Workflow.] Orchestration, revue humaine, export,
        interface.
  \item[Phase 5 --- Graphe (optionnelle).] Extraction d'entités et de
        relations, requêtes multi-sauts.
\end{description}

\subsection{Livrables}

Les livrables attendus sont~: le code source testé, un environnement
reproductible, une documentation technique (interface applicative et
architecture), une documentation utilisateur, un jeu de test (documents et
questions) et les rapports d'évaluation correspondant aux hypothèses H1
à H5.

\subsection{Synthèse de la méthodologie}

La méthodologie se caractérise par une approche modulaire et incrémentale,
une évaluation continue, une validation humaine constitutive, une
traçabilité complète et une adaptabilité volontairement préservée. Les
choix suivants sont particulièrement notables.

\begin{btptable}{Choix méthodologiques et justifications}
\label{tab:choix-methodo}
\begin{tabularx}{\linewidth}{@{}L{4.2cm} Y@{}}
\toprule
\rowcolor{btppale}
\btpth{Choix} & \btpth{Justification} \\
\midrule
Recherche hybride & Meilleur compromis entre termes exacts et concepts. \\
Découpage parent-enfant & Recherche précise et contexte suffisant. \\
Génération par section & Contrôle de qualité, citations, versions. \\
Supervision humaine & Responsabilité, conformité, apprentissage. \\
Workflow déterministe & Traçabilité, coûts bornés, auditabilité. \\
Ingestion asynchrone & L'utilisateur n'attend pas l'indexation. \\
Relationnel et vecteurs & Faits métier d'un côté, similarité de l'autre. \\
\bottomrule
\end{tabularx}
\end{btptable}

Le chapitre suivant incarne ces choix dans une plateforme, en nomme les
composants, et en mesure les performances.

\btpEndMark
